\NSPage{037}
\item[(ii)] The inequalities \NSi{NS-F-P037-001}{a>0}, \NSi{NS-F-P037-002}{P_c>2}, and
\NSi{NS-F-P037-003}{v_s<\mathcal U(P_c,J_c)} hold on \NSi{NS-F-P037-004}{(X_a,X_h]}.
On the first collar \NSi{NS-F-P037-005}{(X_a,X_{\mathrm{an}}]}, one also has
\NSi{NS-F-P037-006}{v_s>2}. At its inner endpoint,
\NSi{NS-F-P037-007}{a(X_a,\eta)>0} and
\NSi{NS-F-P037-008}{v_s(X_a,\eta)>2+c_{\mathrm{ex}}} for a constant
\NSi{NS-F-P037-009}{c_{\mathrm{ex}}>0}. Writing
\NSi{NS-F-P037-010}{y_a=\log(X/X_a)}, the first collar has the factorization
\NSFormula{NS-F-P037-011}{display}{4.33}
\begin{equation}
\begin{gathered}
\mathcal T_0=e^{-t_1^2/y_a^2}g_a(y_a)B_a(y_a,\eta),\qquad g_a(0)>0,\\
B_a(0,\eta)=F(X_a,\eta)\mathfrak s(X_a,\eta)\ne0.
\end{gathered}
\tag{4.33}\label{eq:4.33}
\end{equation}
The functions \NSi{NS-F-P037-012}{g_a,B_a} are smooth through
\NSi{NS-F-P037-013}{y_a=0}, and for every fixed multi-index
\NSi{NS-F-P037-014}{\alpha} there are finite constants with
\NSFormula{NS-F-P037-015}{display}{none}
\[
\lvert\mathcal T_0\rvert\ge c e^{-t_1^2/y_a^2},
\qquad
\lvert\partial_{X,\eta}^{\alpha}\mathcal T_0\rvert
   \le C_{\alpha}e^{-t_1^2/y_a^2}y_a^{-N_{\alpha}}.
\]
All assertions in these two items hold for every
\NSi{NS-F-P037-016}{\eta\in[-1,1]}.

\item[(iii)] Near \NSi{NS-F-P037-017}{X_h}, the fields equal
\NSi{NS-F-P037-018}{U_0=4\eta},
\NSi{NS-F-P037-019}{E_0=P_*f x^{1/10}}, and
\NSi{NS-F-P037-020}{\mathfrak m(X_h,\eta)=\mathfrak m_0(X_h,\eta)}. Here
\NSi{NS-F-P037-021}{\mathfrak m_0=(M_0,I_0,J_0,S_0,C_{p,0})} denotes the five
integrals of this reference pair from zero. These are equalities of functions of
\NSi{NS-F-P037-022}{\eta}, and both pressures use the same prescribed value
\NSi{NS-F-P037-023}{\Pi_0(\eta)} at \NSi{NS-F-P037-024}{X=0}.
\end{enumerate}
Thus \NSi{NS-F-P037-025}{X_R} may exceed any further prescribed lower bound without altering
the earlier choices of \NSi{NS-F-P037-026}{\Lambda,T_{\mathrm{sh}}}.
\end{proposition}

\begin{proof}
We first choose axis data for Proposition~\ref{prop:B.2}, then use
Proposition~\ref{prop:B.3} and Corollary~\ref{cor:B.10} to continue the profile and match its
moments. Set \NSi{NS-F-P037-027}{U_*=4\eta+j_0}, where
\NSi{NS-F-P037-028}{0<j_0\le .05} is the small axis-velocity offset, and define
\NSFormula{NS-F-P037-029}{display}{none}
\[
H_*=D\eta+dU_*,
\qquad
Z_*=-A(1-2\eta U_*)U_*-H_*U_*'-d\Pi_0'+4A\eta\Pi_0.
\]
The pressure bounds in Lemma~\ref{lem:4.8} imply that
\NSi{NS-F-P037-030}{Z_*} is positive at the unique zero of
\NSi{NS-F-P037-031}{H_*}. Choose a threshold
\NSi{NS-F-P037-032}{\delta_*>0}, then a regularization parameter
\NSi{NS-F-P037-033}{\sigma_*>0}, so that
\NSFormula{NS-F-P037-034}{display}{none}
\[
\chi=\frac{H_*^2}{H_*^2+\sigma_*^2}>.99
\quad\hbox{where}\quad \lvert Z_*\rvert\le\delta_*.
\]
These choices specify the positive analytic axis value
\NSFormula{NS-F-P037-035}{display}{none}
\[
\phi(0,\eta)=\exp\!\left(-\Lambda\int_0^{\eta}
\frac{L(w)H_*(w)}{H_*(w)^2+\sigma_*^2}\,dw\right).
\]
The parameter \NSi{NS-F-P037-036}{\Lambda} sets the radial scale
\NSi{NS-F-P037-037}{Y=\Lambda X}. Propositions~\ref{prop:B.2} and~\ref{prop:B.3}
provide the stress-free solution through \NSi{NS-F-P037-038}{Y=4.1} and the exit inequality
\NSi{NS-F-P037-039}{p_{s,1}+p_{s,2}^2/p_{s,1}>2+c_{\mathrm{ex}}} at
\NSi{NS-F-P037-040}{Y=4}. Proposition~\ref{prop:B.5} and Corollary~\ref{cor:B.6}
give the axis and collar regularity in part (i) and the first-collar assertions in part (ii);
\NSi{NS-F-P037-041}{\kappa_0} is the factor reducing the reference shear, and
\NSi{NS-F-P037-042}{t_1} is the width of its initial transition in
\NSi{NS-F-P037-043}{\log(X/X_a)}.

For part (iii), we compute the reference integrals
\NSi{NS-F-P037-044}{\mathfrak m_0}. In the normalization
\NSFormula{NS-F-P037-045}{display}{none}
\[
\widehat{\mathfrak m}=(M/X_R,I/X_R^{3/2},J/X_R^{3/2},S/X_R,C_p),
\]
direct integration of the reference pair gives
\NSFormula{NS-F-P037-046}{display}{4.34}
\begin{equation}
\begin{gathered}
\widehat M_0(x)=4\eta x,\qquad
\widehat I_0(x)=\frac{5\sqrt2}{8}P_*f x^{8/5},\qquad
\widehat J_0(x)=4\eta\widehat I_0(x),\\
\widehat S_0(x)=16\eta^2x-\frac5{12}P_*^2f^2x^{6/5},\qquad
C_{p,0}(x)=\frac52P_*^2f^2x^{1/5}.
\end{gathered}
\tag{4.34}\label{eq:4.34}
\end{equation}
The first continuation supplied by Proposition~\ref{prop:B.5} ends at
\NSFormula{NS-F-P037-047}{display}{none}
\[
X_i=110,\qquad G_i(\eta)=U(X_i,\eta),\qquad
a(X_i,\eta)=4/5,\qquad \mathcal D_XU(X_i,\eta)=0.
\]
The next transition, \eqref{eq:B.34}, holds
\NSi{NS-F-P037-048}{U=G_i} while matching
\NSi{NS-F-P037-049}{E} to the azimuthal reference profile. The value
\NSi{NS-F-P037-050}{U=G_i} is retained until the restoration interval below.

To specify a matching tolerance before \NSi{NS-F-P037-051}{\mathsf C}, put
\NSi{NS-F-P037-052}{\mathcal R=[e^{-8},e^{-5}]}, and take minima and maxima over
\NSi{NS-F-P037-053}{\mathcal R\times[-1,1]}. The reference fields have
\NSi{NS-F-P037-054}{Q_{\min}=\min Q_{s,0}>0},
\NSi{NS-F-P037-055}{e_*=\min E_0>0}, and
\NSPage{038}finite
\NSi{NS-F-P038-001}{w_*=1+\max\lvert N_{s,0}/(E_0Q_{s,0})\rvert}. Here restoration means changing
\NSi{NS-F-P038-002}{U=G_i(\eta)} to \NSi{NS-F-P038-003}{4\eta} on
\NSi{NS-F-P038-004}{e^{-8}<x<e^{-7}}, followed by the five moment corrections on
\NSi{NS-F-P038-005}{e^{-6}<x<e^{-5}}. At the entrance
\NSi{NS-F-P038-006}{x=e^{-6}} to the correction interval, write
\NSi{NS-F-P038-007}{\Delta\widehat{\mathfrak m}=\widehat{\mathfrak m}-\widehat{\mathfrak m}_0},
using \eqref{eq:4.19} with \NSi{NS-F-P038-008}{\rho=X_R}. For the prescribed finite order
\NSi{NS-F-P038-009}{k}, choose \NSi{NS-F-P038-010}{\varepsilon_m>0} so that
\NSFormula{NS-F-P038-011}{display}{none}
\[
\lVert G_i-4\eta\rVert_{C^{k+1}}+
\lVert\Delta\widehat{\mathfrak m}(e^{-6},\mathord\cdot)\rVert_{C^{k+1}}<\varepsilon_m
\]
is sufficient for restoration and correction to satisfy
\NSFormula{NS-F-P038-012}{display}{none}
\[
Q_s\ge Q_{\min}/2,\qquad E\ge e_*/2,\qquad .7\le a\le .9,
\qquad
\left\lvert\frac{N_s}{EQ_s}\right\rvert\le w_*,
\qquad
\lvert b_s\rvert\le\frac{.1}{1+w_*}.
\]
Consequently \NSi{NS-F-P038-013}{v_s=a+b_s^2/a\le .9+.01/.7<1}, while
\NSFormula{NS-F-P038-014}{display}{none}
\[
G=Q_s-\frac{b_sN_s}{aE}\ge\frac67Q_s\ge\frac{3Q_{\min}}7,
\qquad
P_c=\frac{X_Rx}{L}G>2
\]
for every sufficiently large \NSi{NS-F-P038-015}{X_R}. Hence the relaxed condition holds
throughout restoration and moment matching with one fixed tolerance, as required in part (ii).

For each required order \NSi{NS-F-P038-016}{k}, Lemma~\ref{lem:B.7} gives
\NSFormula{NS-F-P038-017}{display}{none}
\[
\lVert G_i-4\eta\rVert_{C^k}
\le C_kj_0+C_k^{\mathrm{nat}}/\Lambda
   +C_k^{\mathrm{join}}(\Lambda)(t_1+\kappa_0+\omega_{\mathrm{fin}}),
\]
where \NSi{NS-F-P038-018}{\omega_{\mathrm{fin}}} is the sum of the two final widths in
\NSi{NS-F-P038-019}{\log X}: the width for setting
\NSi{NS-F-P038-020}{\mathcal D_XU=0} and that for interpolating
\NSi{NS-F-P038-021}{a} to \NSi{NS-F-P038-022}{4/5} before
\NSi{NS-F-P038-023}{X=X_i}. It also bounds
\NSi{NS-F-P038-024}{\log(\mathsf C E(X_i,\mathord\cdot))} independently of the later amplitude, fixing the
length \NSi{NS-F-P038-025}{T_{\mathrm{sh}}} of the transition to
\NSi{NS-F-P038-026}{E_0}. Its endpoint satisfies
\NSFormula{NS-F-P038-027}{display}{none}
\[
x_{\mathrm{sep}}=e^{T_{\mathrm{sh}}}/(\mathsf C P_*)^{10}\longrightarrow0.
\]
Choose \NSi{NS-F-P038-028}{j_0,\Lambda} to control the first two errors, then
\NSi{NS-F-P038-029}{\mathsf C} large, and finally the widths and
\NSi{NS-F-P038-030}{\kappa_0} small. The bound \eqref{eq:B.39} makes the normalized moment
error at \NSi{NS-F-P038-031}{x=e^{-6}} smaller than the prescribed tolerance; the extra
derivative required by \eqref{eq:B.35} is included among the prescribed orders.
Proposition~\ref{prop:B.8} restores \NSi{NS-F-P038-032}{U_0} and solves all five moment
equations before \NSi{NS-F-P038-033}{x=e^{-5}}, preserving the relaxed inequalities above.
Direct integration of \NSi{NS-F-P038-034}{(U_0,E_0)} gives \eqref{eq:4.34}, completing part (iii).
\end{proof}

The inner and outer constructions leave an interval on which only the relaxed cone condition is
known. We will replace its shear by a periodic family with the same mean, every value of which
satisfies the admissible condition. To measure its distance from the four constraint boundaries,
for \NSi{NS-F-P038-035}{a>0}, \NSi{NS-F-P038-036}{b\in\mathbb R}, and
\NSi{NS-F-P038-037}{p=(p_1,p_2)\in\mathbb R^2}, define
\NSFormula{NS-F-P038-038}{display}{4.35}
\begin{equation}
\begin{gathered}
v=a+\frac{b^2}{a},\qquad c=p_1-\frac ba p_2,\qquad j=p_2+\frac ba p_1,\\
\Psi(a,b,p)=\bigl(a,\ v-2,\ c-v,\ 2(c-v)^2-(v-2)j^2\bigr).
\end{gathered}
\tag{4.35}\label{eq:4.35}
\end{equation}
For the profile data \NSi{NS-F-P038-039}{(a,b_s,p_s)}, positivity of all four components is
exactly \NSi{NS-F-P038-040}{a>0}, \NSi{NS-F-P038-041}{v_s>2}, and \eqref{eq:4.22}. The map
\NSi{NS-F-P038-042}{\Psi} is smooth on \NSi{NS-F-P038-043}{\{a>0\}}.

\begin{lemma}\label{lem:4.11}
Let \NSi{NS-F-P038-044}{\mathcal I=[X_-,X_+]\Subset(0,\infty)}, and let
\NSi{NS-F-P038-045}{a,b_s,p_s} be smooth on
\NSi{NS-F-P038-046}{\mathcal I\times[-1,1]}, with
\NSi{NS-F-P038-047}{a>0}. Define \NSi{NS-F-P038-048}{t_s,v_s,P_c,J_c} by
\eqref{eq:4.20}. Suppose \NSi{NS-F-P038-049}{P_c>2} and
\NSi{NS-F-P038-050}{v_s<\mathcal U(P_c,J_c)} throughout this rectangle, and
\NSi{NS-F-P038-051}{v_s>2} in neighborhoods of both radial endpoints. There is a smooth
period-one family \NSi{NS-F-P038-052}{(a_L,-b_L)(X,\eta,\varphi)},
\NSi{NS-F-P038-053}{\varphi\in\mathbb R/\mathbb Z}, such that
\NSFormula{NS-F-P038-054}{display}{none}
\[
\int_0^1(a_L,-b_L)(X,\eta,\varphi)\,d\varphi=(a,-b_s)(X,\eta).
\]
\NSPage{039}Every loop shear satisfies the admissible cone condition with the given vector
\NSi{NS-F-P039-001}{p_s(X,\eta)} held fixed. More precisely, there is
\NSi{NS-F-P039-002}{\mu_L>0} such that
\NSFormula{NS-F-P039-003}{display}{4.36}
\begin{equation}
\Psi_i(a_L,b_L,p_s)\ge\mu_L
\quad\hbox{on }\mathcal I\times[-1,1]\times\mathbb R/\mathbb Z,
\qquad 1\le i\le4.
\tag{4.36}\label{eq:4.36}
\end{equation}
Near both radial endpoints the loop is independent of
\NSi{NS-F-P039-004}{\varphi} and equals
\NSi{NS-F-P039-005}{(a,-b_s)}. Smoothness includes one-sided derivatives at
\NSi{NS-F-P039-006}{\eta=\pm1}; the constant
\NSi{NS-F-P039-007}{\mu_L} may depend on all the fixed input data.
\end{lemma}

\begin{proof}
The variance construction in Lemma~\ref{lem:C.1} supplies a smooth
\NSi{NS-F-P039-008}{2\pi}-periodic ratio
\NSi{NS-F-P039-009}{t_\ell(X,\eta,\theta')} and a smooth
\NSi{NS-F-P039-010}{\rho_\ell(X,\eta)\ge0} such that, writing
\NSi{NS-F-P039-011}{\langle f\rangle=(2\pi)^{-1}\int_0^{2\pi}f\,d\theta'},
\NSFormula{NS-F-P039-012}{display}{none}
\[
\langle t_\ell\rangle=t_s,
\qquad
\langle(t_\ell-t_s)^2\rangle=\rho_\ell/a,
\qquad
v_\ell=v_s+\rho_\ell>2.
\]
It also gives, for every \NSi{NS-F-P039-013}{\theta'},
\NSFormula{NS-F-P039-014}{display}{none}
\[
c_\ell=p_{s,1}+t_\ell p_{s,2}>2,
\qquad
j_\ell=p_{s,2}-t_\ell p_{s,1},
\qquad
v_\ell<\mathcal U(c_\ell,j_\ell).
\]
These are \eqref{eq:C.10} and \eqref{eq:C.8}, together with the cutoff choice
\eqref{eq:C.9}; their construction is smooth also where
\NSi{NS-F-P039-015}{\rho_\ell=0}. Near the radial endpoints it has
\NSi{NS-F-P039-016}{\rho_\ell=0} and \NSi{NS-F-P039-017}{t_\ell=t_s}. Define the lifted
parameter \NSi{NS-F-P039-018}{\varphi}, with
\NSi{NS-F-P039-019}{\varphi(0)=0}, by
\NSFormula{NS-F-P039-020}{display}{4.37}
\begin{equation}
\frac{d\varphi}{d\theta'}=\frac{a(1+t_\ell^2)}{2\pi v_\ell},
\qquad
(a_L,-b_L)=\frac{v_\ell(1,t_\ell)}{1+t_\ell^2}.
\tag{4.37}\label{eq:4.37}
\end{equation}
We verify that the new parameter has period one and that the shear has the prescribed mean in
this parameter. Both follow by direct integration:
\NSFormula{NS-F-P039-021}{display}{none}
\begin{align*}
\int_0^{2\pi}\frac{d\varphi}{d\theta'}\,d\theta'
   &=\frac a{v_\ell}(1+t_s^2+\rho_\ell/a)=1,\\
\int_0^1(a_L,-b_L)\,d\varphi
   &=\frac a{2\pi}\int_0^{2\pi}(1,t_\ell)\,d\theta'
     =(a,at_s)=(a,-b_s).
\end{align*}
The positive derivative in \eqref{eq:4.37} has a positive minimum on the compact parameter set,
so the inverse change of parameter is smooth. The resulting shear satisfies
\NSi{NS-F-P039-022}{-b_L/a_L=t_\ell} and
\NSi{NS-F-P039-023}{a_L(1+t_\ell^2)=v_\ell}. Thus Lemma~\ref{lem:4.5} gives positivity of every
component of \NSi{NS-F-P039-024}{\Psi(a_L,b_L,p_s)}. Their common minimum over
\NSi{NS-F-P039-025}{\mathcal I\times[-1,1]\times\mathbb R/\mathbb Z} is a positive number
\NSi{NS-F-P039-026}{\mu_L}, proving \eqref{eq:4.36}. Where
\NSi{NS-F-P039-027}{\rho_\ell=0}, formula \eqref{eq:4.37} reduces to
\NSi{NS-F-P039-028}{(a_L,-b_L)=(a,-b_s)}, proving the endpoint assertion.
\end{proof}

\subsection{Proof of the leading profile theorem.}\label{sec:4.6}

\begin{proof}[Proof of Theorem~\ref{thm:4.6}]
We first attach the two profiles in Lemma~\ref{lem:4.8} and Proposition~\ref{prop:4.10},
preserving their five cumulative integrals at the joining radius. We then change the shear using
Lemma~\ref{lem:4.11} on the remaining interval where only the relaxed cone holds, estimate the
resulting errors, and use Lemma~\ref{lem:4.7} to solve five equations restoring the integrals. All
estimates below concern \NSi{NS-F-P039-029}{\eta\in[-1,1]}.

\noindent\textbf{Step 1: choose the parameters and join the profiles.}
Choose the outer parameters \NSi{NS-F-P039-030}{M_d,T_d,P_*,\lambda,h} and
\NSi{NS-F-P039-031}{\Pi_0} from Lemma~\ref{lem:4.8}, choosing the family with the endpoint
property in Lemma~\ref{lem:4.9}. Apply Proposition~\ref{prop:4.10} with this same function
\NSi{NS-F-P039-032}{\Pi_0}. In that result \NSi{NS-F-P039-033}{\Lambda} and the logarithmic
transition length \NSi{NS-F-P039-034}{T_{\mathrm{sh}}} are fixed before the amplitude normalization
\NSi{NS-F-P039-035}{\mathsf C}. Let \NSi{NS-F-P039-036}{R_*} be the radius threshold in
 Lemma~\ref{lem:4.8} and \NSi{NS-F-P039-037}{\mathsf C_0} a lower bound from the axis construction after
these choices. We may take
\NSFormula{NS-F-P039-038}{display}{none}
\[
\mathsf C\ge\max\!\left\{\mathsf C_0,\frac{(R_*/110)^{1/10}}{P_*}\right\},
\qquad
X_R=110(\mathsf C P_*)^{10}\ge R_*.
\]
\NSPage{040}The tolerance \NSi{NS-F-P040-001}{\varepsilon_m} and the total logarithmic width
\NSi{NS-F-P040-002}{\omega_{\mathrm{fin}}} are defined in the proof of
Proposition~\ref{prop:4.10}. With the smallness convention in Definition~\ref{def:3.3}, the
profile parameters may be chosen with
\NSFormula{NS-F-P040-003}{display}{none}
\begin{gather*}
0<\mathsf C^{-1}\ll T_{\mathrm{sh}}^{-1}\ll\Lambda^{-1}\ll\sigma_*\ll\delta_*\ll j_0
\ll\varepsilon_m\ll h\ll\lambda\ll P_*^{-1}\ll M_d^{-1}\ll1,\\
0<\omega_{\mathrm{fin}}\ll t_1\ll\kappa_0\ll \mathsf C^{-1}.
\end{gather*}
The second line continues the choices in the first. In particular,
\NSi{NS-F-P040-004}{M_d,P_*} are large,
\NSi{NS-F-P040-005}{\lambda,h,\varepsilon_m,j_0,\delta_*,\sigma_*} are small,
\NSi{NS-F-P040-006}{\Lambda,T_{\mathrm{sh}},\mathsf C} are large, and the three final transition
parameters are small. The exact derived scales and the additional amplitude requirement remain
\NSFormula{NS-F-P040-007}{display}{none}
\[
T_d=e^{M_d}+10,\qquad P_*>e^{T_d},\qquad X_R=110(\mathsf C P_*)^{10}.
\]
After these choices, we fix the periodic shear and the correction bumps in Steps 2--3. We then
choose \NSi{NS-F-P040-008}{N^{-1}\ll\omega_{\mathrm{fin}}}: with the same convention, take
\NSi{NS-F-P040-009}{N} sufficiently large after fixing
\NSi{NS-F-P040-010}{\omega_{\mathrm{fin}}}, those functions, and all constants in the estimates
below. Thus every radial scale and width is fixed before the final integer
\NSi{NS-F-P040-011}{N}.

Put \NSi{NS-F-P040-012}{X_h=X_Re^{-5}}. Write
\NSi{NS-F-P040-013}{(U_{\mathrm{in}},E_{\mathrm{in}})} for the axis connection and
\NSi{NS-F-P040-014}{(U_{\mathrm{out}},E_{\mathrm{out}})} for the prepared outer pair, and define
\NSFormula{NS-F-P040-015}{display}{none}
\[
(U_0,E_0)(X,\eta)=
\begin{cases}
(U_{\mathrm{in}},E_{\mathrm{in}})(X,\eta),&X\le X_h,\\
(U_{\mathrm{out}},E_{\mathrm{out}})(X,\eta),&X\ge X_h.
\end{cases}
\]
Both definitions equal
\NSi{NS-F-P040-016}{(4\eta,P_*(1+\eta^2)^{-1}(X/X_R)^{1/10})} on a neighborhood of
\NSi{NS-F-P040-017}{X_h}, so \NSi{NS-F-P040-018}{U_0,E_0} are smooth there. These are the
profiles before the shear modification in Step 2. Set
\NSi{NS-F-P040-019}{H_0=\sqrt{2X}E_0}, and define
\NSi{NS-F-P040-020}{\mathfrak m_0=(M_0,I_0,J_0,S_0,C_{p,0})} from
\NSi{NS-F-P040-021}{U_0,E_0} by \eqref{eq:4.15}. Retaining the prescribed axis pressure
\NSi{NS-F-P040-022}{\Pi_0(\eta)}, put
\NSi{NS-F-P040-023}{\Pi_0^{\mathrm{prof}}(X,\eta)=\Pi_0(\eta)+C_{p,0}(X,\eta)}. Compute
\NSi{NS-F-P040-024}{V_0,Q_{s,0},N_{s,0},p_{s,0},a_0,b_{s,0},\mathcal T_{0,0}} from
\NSi{NS-F-P040-025}{U_0,E_0,\Pi_0^{\mathrm{prof}}} by \eqref{eq:4.7}, \eqref{eq:4.16}, and
\eqref{eq:4.11}; write \NSi{NS-F-P040-026}{p_0=p_{s,0}},
\NSi{NS-F-P040-027}{b_0=b_{s,0}}, and
\NSi{NS-F-P040-028}{\mathfrak s_0=(a_0,-b_0)}. The five equalities in
Proposition~\ref{prop:4.10}(iii) give
\NSFormula{NS-F-P040-029}{display}{none}
\[
\mathfrak m_0(X_h,\eta)=\mathfrak m_{\mathrm{out}}(X_h,\eta).
\]
For \NSi{NS-F-P040-030}{X\ge X_h} the five integrands also agree, and hence
\NSFormula{NS-F-P040-031}{display}{none}
\begin{align*}
\mathfrak m_0(X,\eta)-\mathfrak m_{\mathrm{out}}(X,\eta)
&=\mathfrak m_0(X_h,\eta)-\mathfrak m_{\mathrm{out}}(X_h,\eta)\\
&\quad+\int_{X_h}^{X}
\begin{pmatrix}
U_0-U_{\mathrm{out}}\\
H_0-H_{\mathrm{out}}\\
U_0H_0-U_{\mathrm{out}}H_{\mathrm{out}}\\
U_0^2-U_{\mathrm{out}}^2-(E_0^2-E_{\mathrm{out}}^2)/2\\
(E_0^2-E_{\mathrm{out}}^2)/(2x)
\end{pmatrix}dx=0.
\end{align*}
By Lemma~\ref{lem:4.4}(i), the pressures and the coefficients
\NSi{NS-F-P040-032}{V_0,Q_{s,0},N_{s,0},p_0,\mathfrak s_0,\mathcal T_{0,0}} agree with the
outer ones on this interval. In particular the pressure normalization and all four total moment
identities in Lemma~\ref{lem:4.8} hold for the joined pair. The angular identity follows from
\NSFormula{NS-F-P040-033}{display}{none}
\[
\int_0^\infty(H_0-H_{\mathrm{out}})\,dX
=I_0(X_h)-I_{\mathrm{out}}(X_h)=0.
\]
The joined pair is regular at the axis, so the conditional stress conclusions of
Lemma~\ref{lem:4.9} now apply. It has zero stress outside
\NSi{NS-F-P040-034}{[X_a,X_b]}, the stated inner and outer edge factorizations, and the
admissible cone near both edges. It satisfies the relaxed cone throughout the intervening interval.

We specify the interval that still needs a shear change. Let
\NSi{NS-F-P040-035}{\mathcal I_1} be the first reserved correction interval from
Lemma~\ref{lem:4.8}, and choose a closed interval
\NSi{NS-F-P040-036}{\mathcal I=[X_-,X_+]\subset(X_a,\inf\mathcal I_1)} containing
\NSPage{041}every point of \NSi{NS-F-P041-001}{(X_a,X_b)} where the joined pair may fail the
admissible condition. Choose its left endpoint strictly between
\NSi{NS-F-P041-002}{X_a} and the end of the analytic collar, and its right endpoint to satisfy
\NSi{NS-F-P041-003}{X_{\mathrm{good}}<X_+<\inf\mathcal I_1}. Shrink the analytic collar
endpoint to a number \NSi{NS-F-P041-004}{X_{\mathrm{an}}\in(X_a,X_-)}. Let
\NSi{NS-F-P041-005}{Y_0<Y_1} lie inside \NSi{NS-F-P041-006}{\mathcal I_1}, with
\NSi{NS-F-P041-007}{X_+<Y_0}, and set
\NSi{NS-F-P041-008}{\mathcal J=[X_-,Y_1]}. The original pair is admissible near the ends of
\NSi{NS-F-P041-009}{\mathcal I} and throughout
\NSi{NS-F-P041-010}{[X_+,Y_1]}. Apply Lemma~\ref{lem:4.11} on
\NSi{NS-F-P041-011}{\mathcal I\times[-1,1]} to its shear
\NSi{NS-F-P041-012}{(a_0,-b_0)} and inviscid vector
\NSi{NS-F-P041-013}{p_0}. For each modified pair
\NSi{NS-F-P041-014}{(U_j,E_j)}, \NSi{NS-F-P041-015}{j=N,f}, below, define
\NSi{NS-F-P041-016}{\mathfrak m_j} by \eqref{eq:4.15}, set
\NSi{NS-F-P041-017}{\Pi_j=\Pi_0+C_{p,j}}, and compute its remaining coefficients by
\eqref{eq:4.7}, \eqref{eq:4.16}, and \eqref{eq:4.11}. We abbreviate
\NSi{NS-F-P041-018}{p_j=p_{s,j}}, \NSi{NS-F-P041-019}{b_j=b_{s,j}}, and write
\NSi{NS-F-P041-020}{V_j} for its radial-velocity profile.

\noindent\textbf{Step 2: realize the periodic shear at a finite frequency.}
The periodic shear has the required cone direction. To obtain it from actual profiles, we integrate
its difference from the original shear and evaluate the resulting antiderivatives at a large radial
frequency. Let \NSi{NS-F-P041-021}{(a_L,-b_L)(X,\eta,\varphi)} be the loop just obtained. Its
mean identities give periodic primitives \NSi{NS-F-P041-022}{\mathscr A,\mathscr B} with zero
mean, satisfying
\NSFormula{NS-F-P041-023}{display}{none}
\[
\partial_\varphi\mathscr A=-\tfrac12(a_L-a_0),
\qquad
\partial_\varphi\mathscr B=\tfrac12E_0(b_L-b_0),
\qquad
\int_0^1\mathscr A\,d\varphi=\int_0^1\mathscr B\,d\varphi=0.
\]
They vanish near both ends of \NSi{NS-F-P041-024}{\mathcal I} because the loop there is
constant. Extend them by zero off \NSi{NS-F-P041-025}{\mathcal I}, and for an integer
\NSi{NS-F-P041-026}{N\ge1} define
\NSFormula{NS-F-P041-027}{display}{4.38}
\begin{equation}
E_N=E_0\exp\!\left(\frac{\mathscr A(X,\eta,N\log X)}{N}\right),
\qquad
U_N=U_0+\frac{\mathscr B(X,\eta,N\log X)}{N}.
\tag{4.38}\label{eq:4.38}
\end{equation}
In the following shear identities, \NSi{NS-F-P041-028}{\mathcal D_X\mathscr A} and
\NSi{NS-F-P041-029}{\mathcal D_X\mathscr B} differentiate \NSi{NS-F-P041-030}{X} with
\NSi{NS-F-P041-031}{\varphi} held fixed; all loop expressions are then evaluated at
\NSi{NS-F-P041-032}{\varphi=N\log X}. The chain rule gives exactly
\NSFormula{NS-F-P041-033}{display}{none}
\begin{align*}
a_N&=1-2\mathcal D_X\log E_N=a_L-\frac{2\mathcal D_X\mathscr A}{N},\\
b_N&=\frac{2\mathcal D_XU_N}{E_N}
=e^{-\mathscr A/N}\left(b_L+\frac{2\mathcal D_X\mathscr B}{NE_0}\right).
\end{align*}
All radii are now fixed and \NSi{NS-F-P041-034}{E_0} has a positive minimum on
\NSi{NS-F-P041-035}{\mathcal J\times[-1,1]}. Since the phase has no
\NSi{NS-F-P041-036}{\eta} dependence, for every fixed
\NSi{NS-F-P041-037}{k} there are constants independent of
\NSi{NS-F-P041-038}{N} such that
\NSFormula{NS-F-P041-039}{display}{4.39}
\begin{equation}
\lVert U_N-U_0,E_N-E_0\rVert_k\le A_k/N,
\qquad
\lVert a_N-a_L,b_N-b_L\rVert_k\le B_k/N.
\tag{4.39}\label{eq:4.39}
\end{equation}
Here and until the end of Step 3, \NSi{NS-F-P041-040}{\lVert\mathord\cdot\rVert_k} takes the
componentwise maximum of \NSi{NS-F-P041-041}{\eta}-derivatives of orders
\NSi{NS-F-P041-042}{0\le j\le k} on
\NSi{NS-F-P041-043}{\mathcal J\times[-1,1]}, or on the indicated subinterval of
\NSi{NS-F-P041-044}{\mathcal J}. The second bound is on
\NSi{NS-F-P041-045}{\mathcal I}; the shear is unchanged on
\NSi{NS-F-P041-046}{\mathcal J\setminus\mathcal I}.

Write \NSi{NS-F-P041-047}{u_N=U_N-U_0},
\NSi{NS-F-P041-048}{e_N=E_N-E_0}, and
\NSi{NS-F-P041-049}{\Delta\mathfrak m_N=\mathfrak m_N-\mathfrak m_0}. There is no change
below \NSi{NS-F-P041-050}{X_-}. For \NSi{NS-F-P041-051}{X\in\mathcal J}, the five exact
increment formulas are
\NSFormula{NS-F-P041-052}{display}{none}
\[
\Delta\mathfrak m_N(X,\eta)=\int_{X_-}^{X}
\begin{pmatrix}
u_N\\
\sqrt{2x}\,e_N\\
H_0u_N+U_0\sqrt{2x}\,e_N+\sqrt{2x}\,u_Ne_N\\
2U_0u_N+u_N^2-E_0e_N-e_N^2/2\\
E_0e_N/x+e_N^2/(2x)
\end{pmatrix}dx.
\]
The interval is fixed and bounded away from zero. The product rule and \eqref{eq:4.39}
therefore give
\NSFormula{NS-F-P041-053}{display}{none}
\[
\lVert\Delta\mathfrak m_N\rVert_k\le D_k/N,
\qquad
\lVert\Pi_N-\Pi_0^{\mathrm{prof}}\rVert_k
=\lVert\Delta C_{p,N}\rVert_k\le D_k/N.
\]
\NSPage{042}Use Lemma~\ref{lem:4.4}(ii) with the common axis value
\NSi{NS-F-P042-001}{\Pi_0(\eta)}, taking \NSi{NS-F-P042-002}{N} large enough that
\NSi{NS-F-P042-003}{E_N\ge\frac12\min_{\mathcal J\times[-1,1]}E_0}. It yields
\NSFormula{NS-F-P042-004}{display}{4.40}
\begin{equation}
\lVert p_N-p_0\rVert_k\le C_k(A_{k+1}+D_{k+1})/N=:P_k/N.
\tag{4.40}\label{eq:4.40}
\end{equation}
Thus the large radial shear change has only a small effect on the integrated vector
\NSi{NS-F-P042-005}{p_s}; the estimate uses no smallness of radial derivatives of
\NSi{NS-F-P042-006}{E_N-E_0} or \NSi{NS-F-P042-007}{U_N-U_0}.

We have controlled the change in \NSi{NS-F-P042-008}{p_s}, while the shear follows the
prescribed loop to order \NSi{NS-F-P042-009}{N^{-1}}. It remains to show that these errors fit
within the loop's strict cone margin. Use the map \NSi{NS-F-P042-010}{\Psi} in
\eqref{eq:4.35}. With \NSi{NS-F-P042-011}{\mathbb T=\mathbb R/\mathbb Z}, let
\NSFormula{NS-F-P042-012}{display}{none}
\[
\mu_L=\min_{\mathcal I\times[-1,1]\times\mathbb T}\min_i
\Psi_i(a_L,b_L,p_0)>0,
\qquad
\mu_R=\min_{[X_+,Y_1]\times[-1,1]}\min_i
\Psi_i(a_0,b_0,p_0)>0.
\]
Both minima are positive by Lemma~\ref{lem:4.11} and the choice of
\NSi{NS-F-P042-013}{\mathcal I}. Fix a compact neighborhood of the values of
\NSi{NS-F-P042-014}{(a_L,b_L,p_0)} and \NSi{NS-F-P042-015}{(a_0,b_0,p_0)} over the two
displayed domains, on which \NSi{NS-F-P042-016}{a>0}, and let
\NSi{NS-F-P042-017}{C_\Psi} be an upper bound on the Lipschitz constant of
\NSi{NS-F-P042-018}{\Psi} on that neighborhood. By \eqref{eq:4.39} and \eqref{eq:4.40}, for
all sufficiently large \NSi{NS-F-P042-019}{N} the perturbed triples lie in that neighborhood and
satisfy, componentwise,
\NSFormula{NS-F-P042-020}{display}{4.41}
\begin{equation}
\begin{aligned}
\Psi_i(a_N,b_N,p_N)&\ge\mu_L-C_\Psi(B_0+P_0)/N\ge\mu_L/2
&&\text{on }\mathcal I,\\
\Psi_i(a_N,b_N,p_N)&\ge\mu_R-C_\Psi P_0/N\ge\mu_R/2
&&\text{on }[X_+,Y_1].
\end{aligned}
\tag{4.41}\label{eq:4.41}
\end{equation}
These four positive quantities assert \NSi{NS-F-P042-021}{a_N>0},
\NSi{NS-F-P042-022}{v_N>2}, and \eqref{eq:4.22}. The modulated pair is therefore admissible
throughout \NSi{NS-F-P042-023}{\mathcal J}. Its five cumulative integrals still differ from the
original ones by \NSi{NS-F-P042-024}{O(N^{-1})}; we now restore them exactly.

\noindent\textbf{Step 3: solve the five moment equations.}
On \NSi{NS-F-P042-025}{(Y_0,Y_1)\subset\mathcal I_1}, modulation has not changed the profiles,
and
\NSFormula{NS-F-P042-026}{display}{none}
\[
U_N=U_0=0,
\qquad
E_N=E_0=K(\eta)X^{-1/2-\lambda},
\qquad
K(\eta)=c_{\mathrm{patch}}(1+\eta^2)^{-1}>0.
\]
Choose two nonnegative nonzero smooth functions
\NSi{NS-F-P042-027}{\beta_1,\beta_2} with ordered disjoint compact supports in
\NSi{NS-F-P042-028}{(Y_0,Y_1)}, and three such functions
\NSi{NS-F-P042-029}{\gamma_1,\gamma_2,\gamma_3}. Set
\NSFormula{NS-F-P042-030}{display}{none}
\[
u_c=\sum_{i=1}^2\alpha_i(\eta)\beta_i(X),
\qquad
e_c=\sum_{j=1}^3\xi_j(\eta)\gamma_j(X),
\qquad
c=(\alpha_1,\alpha_2,\xi_1,\xi_2,\xi_3).
\]
For \NSi{NS-F-P042-031}{U_f=U_N+u_c},
\NSi{NS-F-P042-032}{E_f=E_N+e_c}, the exact changes at
\NSi{NS-F-P042-033}{Y_1}, ordered as
\NSi{NS-F-P042-034}{(M,J,I,S,C_p)}, are
\NSFormula{NS-F-P042-035}{display}{4.42}
\begin{equation}
\begin{pmatrix}
\int u_c\,dX\\
\int(H_0u_c+\sqrt{2X}\,u_ce_c)\,dX\\
\int\sqrt{2X}\,e_c\,dX\\
\int(u_c^2-E_0e_c-e_c^2/2)\,dX\\
\int(E_0e_c/X+e_c^2/(2X))\,dX
\end{pmatrix}
=B(\eta)c+\mathcal Q_\eta(c,c)=d_N(\eta).
\tag{4.42}\label{eq:4.42}
\end{equation}
All five integrals are over \NSi{NS-F-P042-036}{(Y_0,Y_1)}, and
\NSi{NS-F-P042-037}{d_N} is
\NSi{NS-F-P042-038}{-\Delta\mathfrak m_N(Y_1,\eta)} in the same row order. The linear map has
two diagonal blocks:
\NSFormula{NS-F-P042-039}{display}{none}
\[
B_U=
\begin{pmatrix}
\int\beta_1&\int\beta_2\\
\int H_0\beta_1&\int H_0\beta_2
\end{pmatrix},
\qquad
B_E=
\begin{pmatrix}
\int\sqrt{2X}\gamma_1&\int\sqrt{2X}\gamma_2&\int\sqrt{2X}\gamma_3\\
-\int E_0\gamma_1&-\int E_0\gamma_2&-\int E_0\gamma_3\\
\int E_0\gamma_1/X&\int E_0\gamma_2/X&\int E_0\gamma_3/X
\end{pmatrix}.
\]
The weights in \NSi{NS-F-P042-040}{B_U} have exponents
\NSi{NS-F-P042-041}{0,-\lambda}, and those in
\NSi{NS-F-P042-042}{B_E} have exponents
\NSi{NS-F-P042-043}{1/2,-1/2-\lambda,-3/2-\lambda}. They are distinct since
\NSi{NS-F-P042-044}{\lambda>0}. By Lemma~\ref{lem:4.7}, both blocks are invertible,
